\begin{abstract}
Model merging has emerged as a powerful paradigm for combining multiple task-specific expert neural networks without the prohibitive computational cost of joint multi-task retraining. 
However, existing methods—including state-of-the-art representation-alignment and test-time adaptive techniques—rely fundamentally on linear combinations or rigid, predefined manifold projections of parameter vectors. 
These methods are inherently limited by the highly non-convex, curved, and disjoint geometry of neural loss landscapes, where straight-line paths between expert parameters can cross high-loss barriers, degrading multi-task capability. 
In this work, we explore this flat-space paradigm and introduce \textbf{FoldMerge (Neural Origami)}, a non-linear coordinate-transformation framework that treats model merging as a continuous weight-space warping process. 
Instead of averaging parameters in Euclidean weight-space, FoldMerge investigates the use of a differentiable weight-space diffeomorphism parameterized by normalizing flows to map disjoint parameter basins into a latent shared coordinate system (``Origami Space''). 
We compute the additive latent combination of the warped parameters in Origami Space and project the merged coordinates back via the analytical inverse diffeomorphism. 
To ensure smooth, structure-preserving transformations with zero extra computational overhead, we introduce an implicit flow regularization penalty via parameter-wise $\ell_2$ weight decay on the flow parameters. 
Evaluated on the 8-task Vision-Language ViT-B/32 benchmark, FoldMerge provides a compelling proof-of-concept, achieving an average accuracy of \textbf{89.76\%} on par with the highly optimized SyMerge framework (89.74\%). 
Our work demonstrates that non-linear parameter-space warping is computationally trainable and viable, and we discuss its unique limitations (including coordinate-dependence and parameter overhead) to open up fresh directions for future deep weight-space geometry research.
\end{abstract}
